% Options for packages loaded elsewhere
\PassOptionsToPackage{unicode}{hyperref}
\PassOptionsToPackage{hyphens}{url}
\documentclass[
]{article}
\usepackage{xcolor}
\usepackage{amsmath,amssymb}
\setcounter{secnumdepth}{-\maxdimen} % remove section numbering
\usepackage{iftex}
\ifPDFTeX
  \usepackage[T1]{fontenc}
  \usepackage[utf8]{inputenc}
  \usepackage{textcomp} % provide euro and other symbols
\else % if luatex or xetex
  \usepackage{unicode-math} % this also loads fontspec
  \defaultfontfeatures{Scale=MatchLowercase}
  \defaultfontfeatures[\rmfamily]{Ligatures=TeX,Scale=1}
\fi
\usepackage{lmodern}
\ifPDFTeX\else
  % xetex/luatex font selection
\fi
% Use upquote if available, for straight quotes in verbatim environments
\IfFileExists{upquote.sty}{\usepackage{upquote}}{}
\IfFileExists{microtype.sty}{% use microtype if available
  \usepackage[]{microtype}
  \UseMicrotypeSet[protrusion]{basicmath} % disable protrusion for tt fonts
}{}
\makeatletter
\@ifundefined{KOMAClassName}{% if non-KOMA class
  \IfFileExists{parskip.sty}{%
    \usepackage{parskip}
  }{% else
    \setlength{\parindent}{0pt}
    \setlength{\parskip}{6pt plus 2pt minus 1pt}}
}{% if KOMA class
  \KOMAoptions{parskip=half}}
\makeatother
\usepackage{longtable,booktabs,array}
\usepackage{calc} % for calculating minipage widths
% Correct order of tables after \paragraph or \subparagraph
\usepackage{etoolbox}
\makeatletter
\patchcmd\longtable{\par}{\if@noskipsec\mbox{}\fi\par}{}{}
\makeatother
% Allow footnotes in longtable head/foot
\IfFileExists{footnotehyper.sty}{\usepackage{footnotehyper}}{\usepackage{footnote}}
\makesavenoteenv{longtable}
\usepackage{graphicx}
\makeatletter
\newsavebox\pandoc@box
\newcommand*\pandocbounded[1]{% scales image to fit in text height/width
  \sbox\pandoc@box{#1}%
  \Gscale@div\@tempa{\textheight}{\dimexpr\ht\pandoc@box+\dp\pandoc@box\relax}%
  \Gscale@div\@tempb{\linewidth}{\wd\pandoc@box}%
  \ifdim\@tempb\p@<\@tempa\p@\let\@tempa\@tempb\fi% select the smaller of both
  \ifdim\@tempa\p@<\p@\scalebox{\@tempa}{\usebox\pandoc@box}%
  \else\usebox{\pandoc@box}%
  \fi%
}
% Set default figure placement to htbp
\def\fps@figure{htbp}
\makeatother
\setlength{\emergencystretch}{3em} % prevent overfull lines
\providecommand{\tightlist}{%
  \setlength{\itemsep}{0pt}\setlength{\parskip}{0pt}}
\usepackage{bookmark}
\IfFileExists{xurl.sty}{\usepackage{xurl}}{} % add URL line breaks if available
\urlstyle{same}
\hypersetup{
  hidelinks,
  pdfcreator={LaTeX via pandoc}}

\author{}
\date{}

\begin{document}

\subsubsection{МИНИСТЕРСТВО НАУКИ И ВЫСШЕГО ОБРАЗОВАНИЯ РОССИЙСКОЙ
ФЕДЕРАЦИИ}\label{ux43cux438ux43dux438ux441ux442ux435ux440ux441ux442ux432ux43e-ux43dux430ux443ux43aux438-ux438-ux432ux44bux441ux448ux435ux433ux43e-ux43eux431ux440ux430ux437ux43eux432ux430ux43dux438ux44f-ux440ux43eux441ux441ux438ux439ux441ux43aux43eux439-ux444ux435ux434ux435ux440ux430ux446ux438ux438}

\subsubsection{Федерального государственного автономного
образовательного учреждения высшего
образования}\label{ux444ux435ux434ux435ux440ux430ux43bux44cux43dux43eux433ux43e-ux433ux43eux441ux443ux434ux430ux440ux441ux442ux432ux435ux43dux43dux43eux433ux43e-ux430ux432ux442ux43eux43dux43eux43cux43dux43eux433ux43e-ux43eux431ux440ux430ux437ux43eux432ux430ux442ux435ux43bux44cux43dux43eux433ux43e-ux443ux447ux440ux435ux436ux434ux435ux43dux438ux44f-ux432ux44bux441ux448ux435ux433ux43e-ux43eux431ux440ux430ux437ux43eux432ux430ux43dux438ux44f}

\subsubsection{«НАЦИОНАЛЬНЫЙ ИССЛЕДОВАТЕЛЬСКИЙ ТЕХНОЛОГИЧЕСКИЙ
УНИВЕРСИТЕТ
МИСИС»}\label{ux43dux430ux446ux438ux43eux43dux430ux43bux44cux43dux44bux439-ux438ux441ux441ux43bux435ux434ux43eux432ux430ux442ux435ux43bux44cux441ux43aux438ux439-ux442ux435ux445ux43dux43eux43bux43eux433ux438ux447ux435ux441ux43aux438ux439-ux443ux43dux438ux432ux435ux440ux441ux438ux442ux435ux442-ux43cux438ux441ux438ux441}

\subsubsection{Институт компьютерных наук Кафедра Инженерной
кибернетики}\label{ux438ux43dux441ux442ux438ux442ux443ux442-ux43aux43eux43cux43fux44cux44eux442ux435ux440ux43dux44bux445-ux43dux430ux443ux43a-ux43aux430ux444ux435ux434ux440ux430-ux438ux43dux436ux435ux43dux435ux440ux43dux43eux439-ux43aux438ux431ux435ux440ux43dux435ux442ux438ux43aux438}

\subsubsection{Лабораторная работа №1 по дисциплине «Современные
интеллектуальные сетевые
сервисы»}\label{ux43bux430ux431ux43eux440ux430ux442ux43eux440ux43dux430ux44f-ux440ux430ux431ux43eux442ux430-1-ux43fux43e-ux434ux438ux441ux446ux438ux43fux43bux438ux43dux435-ux441ux43eux432ux440ux435ux43cux435ux43dux43dux44bux435-ux438ux43dux442ux435ux43bux43bux435ux43aux442ux443ux430ux43bux44cux43dux44bux435-ux441ux435ux442ux435ux432ux44bux435-ux441ux435ux440ux432ux438ux441ux44b}

Проектирование AI-сервиса на FastAPI

\begin{quote}
\textbf{Выполнил:} Студент группы МПИ-25-1 Харлашкина А.В.

\textbf{Проверил:} Преподаватель Корнет М.Е.
\end{quote}

\subsubsection{\texorpdfstring{Оценка:
}{Оценка: }}\label{ux43eux446ux435ux43dux43aux430}

\begin{quote}
Москва, 2026
\end{quote}

СОДЕРЖАНИЕ

\hyperref[ux446ux435ux43bux44c-ux440ux430ux431ux43eux442ux44b]{1 Цель
работы
\hyperref[ux446ux435ux43bux44c-ux440ux430ux431ux43eux442ux44b]{3}}

\hyperref[ux437ux430ux434ux430ux447ux438-ux440ux430ux431ux43eux442ux44b]{2
Задачи работы
\hyperref[ux437ux430ux434ux430ux447ux438-ux440ux430ux431ux43eux442ux44b]{3}}

\hyperref[ux438ux441ux445ux43eux434ux43dux44bux435-ux434ux430ux43dux43dux44bux435]{3
Исходные данные
\hyperref[ux438ux441ux445ux43eux434ux43dux44bux435-ux434ux430ux43dux43dux44bux435]{3}}

\hyperref[ux442ux440ux435ux431ux443ux435ux43cux44bux435-endpoints]{4
Требуемые endpoints
\hyperref[ux442ux440ux435ux431ux443ux435ux43cux44bux435-endpoints]{4}}

\hyperref[ux432ux44bux43fux43eux43bux43dux435ux43dux438ux435-ux440ux430ux431ux43eux442ux44b]{5
Выполнение работы
\hyperref[ux432ux44bux43fux43eux43bux43dux435ux43dux438ux435-ux440ux430ux431ux43eux442ux44b]{5}}

\hyperref[ux441ux43eux437ux434ux430ux43dux438ux435-ux43fux440ux43eux435ux43aux442ux430]{5.1
Создание проекта
\hyperref[ux441ux43eux437ux434ux430ux43dux438ux435-ux43fux440ux43eux435ux43aux442ux430]{5}}

\hyperref[ux441ux43eux437ux434ux430ux43dux438ux435-fastapi-ux43fux440ux438ux43bux43eux436ux435ux43dux438ux44f]{5.2
Создание FastAPI-приложения
\hyperref[ux441ux43eux437ux434ux430ux43dux438ux435-fastapi-ux43fux440ux438ux43bux43eux436ux435ux43dux438ux44f]{5}}

\hyperref[ux43fux435ux440ux432ux44bux439-endpoint-health]{5.2.1 Первый
endpoint -- /health
\hyperref[ux43fux435ux440ux432ux44bux439-endpoint-health]{6}}

\hyperref[ux434ux43eux431ux430ux432ux43bux435ux43dux438ux435-ux438ux43dux444ux43eux440ux43cux430ux446ux438ux438-ux43e-ux43cux43eux434ux435ux43bux438-endpoint-model-info]{5.2.2
Добавление информации о модели: endpoint /model-info
\hyperref[ux434ux43eux431ux430ux432ux43bux435ux43dux438ux435-ux438ux43dux444ux43eux440ux43cux430ux446ux438ux438-ux43e-ux43cux43eux434ux435ux43bux438-endpoint-model-info]{7}}

\hyperref[ux43fux440ux43eux435ux43aux442ux438ux440ux43eux432ux430ux43dux438ux435-ux432ux445ux43eux434ux43dux44bux445-ux434ux430ux43dux43dux44bux445]{5.2.3
Проектирование входных данных
\hyperref[ux43fux440ux43eux435ux43aux442ux438ux440ux43eux432ux430ux43dux438ux435-ux432ux445ux43eux434ux43dux44bux445-ux434ux430ux43dux43dux44bux445]{8}}

\hyperref[ux440ux435ux430ux43bux438ux437ux430ux446ux438ux44f-ux43cux43eux434ux435ux43bux438]{5.2.4
Реализация модели
\hyperref[ux440ux435ux430ux43bux438ux437ux430ux446ux438ux44f-ux43cux43eux434ux435ux43bux438]{9}}

\hyperref[ux440ux435ux430ux43bux438ux437ux430ux446ux438ux44f-post-predict]{5.2.5
Реализация POST /predict
\hyperref[ux440ux435ux430ux43bux438ux437ux430ux446ux438ux44f-post-predict]{10}}

\hyperref[ux43eux440ux433ux430ux43dux438ux437ux430ux446ux438ux44f-ux445ux440ux430ux43dux435ux43dux438ux44f-ux440ux435ux437ux443ux43bux44cux442ux430ux442ux43eux432]{5.2.6
Организация хранения результатов
\hyperref[ux43eux440ux433ux430ux43dux438ux437ux430ux446ux438ux44f-ux445ux440ux430ux43dux435ux43dux438ux44f-ux440ux435ux437ux443ux43bux44cux442ux430ux442ux43eux432]{11}}

\hyperref[ux440ux435ux430ux43bux438ux437ux430ux446ux438ux44f-endpoint-predictionsrequest_id]{5.2.7
Реализация endpoint /predictions/\{request\_id\}
\hyperref[ux440ux435ux430ux43bux438ux437ux430ux446ux438ux44f-endpoint-predictionsrequest_id]{12}}

\hyperref[ux440ux435ux430ux43bux438ux437ux430ux446ux438ux44f-endpoint-predictions]{5.2.8
Реализация endpoint /predictions
\hyperref[ux440ux435ux430ux43bux438ux437ux430ux446ux438ux44f-endpoint-predictions]{14}}

\hyperref[ux43fux440ux43eux432ux435ux440ux43aux430-ux431ux438ux437ux43dux435ux441-ux43eux433ux440ux430ux43dux438ux447ux435ux43dux438ux439]{5.2.9
Проверка бизнес ограничений
\hyperref[ux43fux440ux43eux432ux435ux440ux43aux430-ux431ux438ux437ux43dux435ux441-ux43eux433ux440ux430ux43dux438ux447ux435ux43dux438ux439]{16}}

\hyperref[openapi-ux43aux430ux43a-ux43aux43eux43dux442ux440ux430ux43aux442]{5.3
OpenAPI как контракт
\hyperref[openapi-ux43aux430ux43a-ux43aux43eux43dux442ux440ux430ux43aux442]{18}}

\hyperref[ux442ux435ux441ux442ux438ux440ux43eux432ux430ux43dux438ux435-api]{5.4
Тестирование API
\hyperref[ux442ux435ux441ux442ux438ux440ux43eux432ux430ux43dux438ux435-api]{18}}

\hyperref[ux430ux43dux430ux43bux438ux437-api-ux43aux430ux43a-ux43aux43eux43dux442ux440ux430ux43aux442ux430]{5.5
Анализ API как контракта
\hyperref[ux430ux43dux430ux43bux438ux437-api-ux43aux430ux43a-ux43aux43eux43dux442ux440ux430ux43aux442ux430]{21}}

\hyperref[ux440ux430ux437ux434ux435ux43bux435ux43dux438ux435-ux43fux440ux438ux43bux43eux436ux435ux43dux438ux44f-ux43dux430-ux43cux43eux434ux443ux43bux438-ux430ux440ux445ux438ux442ux435ux43aux442ux443ux440ux430-ux43fux440ux438ux43bux43eux436ux435ux43dux438ux44f]{5.6
Разделение приложения на модули: архитектура приложения
\hyperref[ux440ux430ux437ux434ux435ux43bux435ux43dux438ux435-ux43fux440ux438ux43bux43eux436ux435ux43dux438ux44f-ux43dux430-ux43cux43eux434ux443ux43bux438-ux430ux440ux445ux438ux442ux435ux43aux442ux443ux440ux430-ux43fux440ux438ux43bux43eux436ux435ux43dux438ux44f]{21}}

\hyperref[ux432ux44bux432ux43eux434]{ВЫВОД
\hyperref[ux432ux44bux432ux43eux434]{23}}

\hyperref[ux43fux440ux438ux43bux43eux436ux435ux43dux438ux435-ux430-ux43fux43eux43bux43dux44bux439-ux43aux43eux434-ux43fux440ux438ux43bux43eux436ux435ux43dux438ux44f]{Приложение
А: Полный код приложения
\hyperref[ux43fux440ux438ux43bux43eux436ux435ux43dux438ux435-ux430-ux43fux43eux43bux43dux44bux439-ux43aux43eux434-ux43fux440ux438ux43bux43eux436ux435ux43dux438ux44f]{24}}

\section{Цель
работы}\label{ux446ux435ux43bux44c-ux440ux430ux431ux43eux442ux44b}

Цель практической работы -- научиться проектировать программную
архитектуру AI-сервиса, использующего модель машинного обучения,
определить место API в общей ИИ-системе и реализовать REST API на
FastAPI, изучив основные принципы программного взаимодействия между
AI-моделью и внешними информационными системами.

\section{Задачи
работы}\label{ux437ux430ux434ux430ux447ux438-ux440ux430ux431ux43eux442ux44b}

В ходе работы необходимо:

\begin{itemize}
\item
  определить потребителей AI-сервиса;
\item
  определить формат входных и выходных данных;
\item
  спроектировать контур инференса;
\item
  разделить синхронные и асинхронные операции;
\item
  определить основные программные компоненты системы;
\item
  выбрать протоколы взаимодействия;
\item
  сформировать структурированный контракт REST API для AI-сервиса;
\item
  определить способы хранения результатов;
\item
  определить требования к логированию и наблюдаемости;
\item
  выявить основные архитектурные риски;
\item
  создать приложение на FastAPI и реализовать REST endpoints;
\item
  изучить методы GET и POST, передачу данных через JSON, использование
  path- и query-параметров;
\item
  описывать входные и выходные схемы с помощью Pydantic и выполнять
  валидацию входных данных;
\item
  возвращать корректные HTTP-коды и обрабатывать ошибки;
\item
  изучить автоматически создаваемую OpenAPI-документацию и
  протестировать API через Swagger UI;
\item
  реализовать минимальный прототип API на FastAPI.
\end{itemize}

\section{Исходные
данные}\label{ux438ux441ux445ux43eux434ux43dux44bux435-ux434ux430ux43dux43dux44bux435}

Банк разрабатывает сервис агроскоринга. Система должна оценивать риск
при работе с сельскохозяйственными предприятиями. Пользователь или
другая информационная система передаёт сведения о хозяйстве:

\begin{itemize}
\item
  Регион
\item
  площадь посевов;
\item
  тип основной культуры;
\item
  погодные показатели;
\item
  историю платежей;
\item
  количество просроченных платежей;
\item
  размер текущей задолженности.
\end{itemize}

На основании этих данных модель машинного обучения должна определить
риск профиля предприятия. Результат должен содержать:

\begin{itemize}
\item
  числовую оценку риска;
\item
  категорию риска;
\item
  рекомендацию оператору банка.
\end{itemize}

\begin{quote}
\{

" risk\_score ": 0.73 , " risk\_level": " high ", " recommendation ": ""

\}
\end{quote}

Пример ответа:

\begin{quote}
1

2

3

4

5

Листинг 1: Пример ответа /predict
\end{quote}

Сервис должен позволять:

\begin{itemize}
\item
  проверить, работает ли API;
\item
  получить информацию о модели;
\item
  выполнить оценку риска одного хозяйства;
\item
  получить сохранённый результат прогноза;
\item
  получить несколько последних прогнозов;
\item
  корректно сообщать клиенту об ошибках.
\end{itemize}

Предполагается, что в будущем этим API будут пользоваться: CRM банка,
веб-портал, мобильное приложение, внутренние банковские сервисы,
BI-система.

\section{Требуемые
endpoints}\label{ux442ux440ux435ux431ux443ux435ux43cux44bux435-endpoints}

API должен иметь следующие endpoints:

Таблица №1 -- Список endpoints API

\begin{longtable}[]{@{}
  >{\raggedright\arraybackslash}p{(\linewidth - 4\tabcolsep) * \real{0.2297}}
  >{\raggedright\arraybackslash}p{(\linewidth - 4\tabcolsep) * \real{0.3701}}
  >{\raggedright\arraybackslash}p{(\linewidth - 4\tabcolsep) * \real{0.4001}}@{}}
\toprule\noalign{}
\begin{minipage}[b]{\linewidth}\centering
\textbf{Метод}
\end{minipage} & \begin{minipage}[b]{\linewidth}\centering
\textbf{Endpoint}
\end{minipage} & \begin{minipage}[b]{\linewidth}\centering
\textbf{Назначение}
\end{minipage} \\
\midrule\noalign{}
\endhead
\bottomrule\noalign{}
\endlastfoot
GET & /health & Проверка работоспособности API \\
GET & /model-info & Информация о модели \\
POST & /predict & Оценка риска хозяйства \\
GET & /predictions/\{request\_id\} & Получение прогноза по ID \\
GET & /predictions & Список прогнозов \\
\end{longtable}

\section{Выполнение
работы}\label{ux432ux44bux43fux43eux43bux43dux435ux43dux438ux435-ux440ux430ux431ux43eux442ux44b}

\section{Создание
проекта}\label{ux441ux43eux437ux434ux430ux43dux438ux435-ux43fux440ux43eux435ux43aux442ux430}

Создадим каталог agro\_api/ (Рисунок 1). Также создадим виртуальное
окружение и установим необходимые зависимости.

\begin{figure}
\centering
\includegraphics[width=3.16948in,height=2.63575in]{media/image1.png}
\caption{Рисунок 1 -- Структура проекта agro\_api}
\end{figure}

\section{Создание
FastAPI-приложения}\label{ux441ux43eux437ux434ux430ux43dux438ux435-fastapi-ux43fux440ux438ux43bux43eux436ux435ux43dux438ux44f}

Создадим файл main.py и добавим в него объект app = FastAPI(), который
представляет наше API-приложение.

\begin{figure}
\centering
\includegraphics[width=6.46606in,height=1.31714in]{media/image2.png}
\caption{Рисунок 2 -- Создание объекта app}
\end{figure}

\section{Первый endpoint --
/health}\label{ux43fux435ux440ux432ux44bux439-endpoint-health}

Создадим Endpoint /health (Рисунок 3).

\includegraphics[width=6.26328in,height=1.48235in]{media/image3.png}

Рисунок 3 -- Создание Endpoint /health

Откроем http://127.0.0.1:8000/health. Получили результат: \{" status ":
"ok"\} (Рисунок 4).

\begin{figure}
\centering
\includegraphics[width=4.08675in,height=1.87936in]{media/image4.png}
\caption{Рисунок 4 -- Результат /health}
\end{figure}

Теперь откроем http://127.0.0.1:8000/docs. FastAPI автоматически создаёт
интерактивную документацию API. В ней появился GET /health (Рисунок 5).

Также FastAPI автоматически формирует описание API по стандарту OpenAPI.
Откроем http: //127.0.0.1:8000/openapi.json и увидим JSON-документ,
описывающий API (Рисунок 6).

\begin{figure}
\centering
\includegraphics[width=6.9983in,height=6.43936in]{media/image5.png}
\caption{Рисунок 5 -- Результат в Swagger UI}
\end{figure}

\begin{figure}
\centering
\includegraphics[width=6.41545in,height=1.90501in]{media/image6.png}
\caption{Рисунок 6 -- Фрагмент OpenAPI-спецификации}
\end{figure}

\section{Добавление информации о модели: endpoint
/model-info}\label{ux434ux43eux431ux430ux432ux43bux435ux43dux438ux435-ux438ux43dux444ux43eux440ux43cux430ux446ux438ux438-ux43e-ux43cux43eux434ux435ux43bux438-endpoint-model-info}

Теперь создадим endpoint GET /model-info (Рисунок 7).

\begin{figure}
\centering
\includegraphics[width=6.61895in,height=2.60987in]{media/image7.png}
\caption{Рисунок 7 -- Создание Endpoint /model-info}
\end{figure}

Ответ на впорос чем /health отличается от /model-info?

/health -- техническая проверка доступности API: сервис отвечает
\{"status": "ok"\}.

/model-info -- описание используемой модели: имя, версия, тип и
состояние готовности. То есть health характеризует работоспособность
сервиса, а model-info -- состояние и метаданные модели.

\section{Проектирование входных
данных}\label{ux43fux440ux43eux435ux43aux442ux438ux440ux43eux432ux430ux43dux438ux435-ux432ux445ux43eux434ux43dux44bux445-ux434ux430ux43dux43dux44bux445}

Для выполнения прогноза клиент должен передать сведения о хозяйстве.
FastAPI должен проверить эти данные до передачи модели. FastAPI
использует Pydantic для описания структуры данных. Поэтому добавим
импорт Pydantic и создадим модель FarmRequest (Рисунок 8).

\begin{figure}
\centering
\includegraphics[width=6.03104in,height=1.95824in]{media/image8.png}
\caption{Рисунок 8 -- Создание модели FarmRequest}
\end{figure}

Пояснения: area\_ha: float = Field(gt=0) означает тип данных float,
значение строго больше 0. Если же клиент передаст: area\_ha = -200, то
FastAPI обнаружит ошибку ещё до запуска функции инференса.

\section{Реализация
модели}\label{ux440ux435ux430ux43bux438ux437ux430ux446ux438ux44f-ux43cux43eux434ux435ux43bux438}

В данной работе вместо модели машинного обучения будем использовать
простую функцию, реализующую расчёт риска по некоторым финансовым и
отраслевым показателям. Добавим функцию расчета риска (Рисунок 9).

\begin{figure}
\centering
\includegraphics[width=4.41705in,height=3.25109in]{media/image9.png}
\caption{Рисунок 9 -- Функция расчета риска}
\end{figure}

Также создадим функции risk\_level (Рисунок 10) и recommendation
(Рисунок 11):

\begin{itemize}
\item
  risk\_level преобразует числовой risk\_score в понятную категорию
  риска: low, medium или high.
\item
  recommendation формирует рекомендацию для банка в зависимости от
  полученного уровня риска.
\end{itemize}

\begin{figure}
\centering
\includegraphics[width=5.1005in,height=1.31694in]{media/image10.png}
\caption{Рисунок 10 -- Функция risk\_level}
\end{figure}

\begin{figure}
\centering
\includegraphics[width=5.1005in,height=1.27761in]{media/image10.png}
\caption{Рисунок 11 -- Функция recommendation}
\end{figure}

\section{Реализация POST
/predict}\label{ux440ux435ux430ux43bux438ux437ux430ux446ux438ux44f-post-predict}

Создадим Endpoint POST /predict (Рисунок 12). Для этого импортуруем
uuid. Данный endpoint является основным endpoint разработанного
AI-сервиса, поскольку именно через него передаются данные о хозяйстве
для выполнения оценки кредитного риска.

Также реализуем добавление кода ответа в POST /predict. Какой статус
корректнее для инференса --- «200 OK» или «201 Created»? В данной
реализации выбран 200 OK. Endpoint выполняет инференс и возвращает
результат операции; отдельный ресурс с самостоятельным жизненным циклом
здесь не создаётся. Важно последовательно использовать выбранную
семантику.

Для возвращаемого результата создадим PredictionResponse, так как
внешняя система должна получать ожидаемую структуру данных независимо от
внут-реннего устройства модели.

\begin{figure}
\centering
\includegraphics[width=6.23501in,height=4.80675in]{media/image11.png}
\caption{Рисунок 12 -- Endpoint POST /predict}
\end{figure}

Такое разделение позволяет отделить HTTP-интерфейс от бизнес-логики и
расчёта риска, что упрощает тестирование и дальнейшую замену
демонстрационной модели на полноценную ML-модель.

Проверим через Swagger: откроем /docs, выберем POST /predict и введем
тестовый запрос (Рисунок 13).

\begin{figure}
\centering
\includegraphics[width=3.22276in,height=2.49526in]{media/image12.png}
\caption{Рисунок 13 -- Тестовый запрос}
\end{figure}

После чего получим результат (Рисунок 14).

\begin{figure}
\centering
\includegraphics[width=6.65258in,height=4.72225in]{media/image13.png}
\caption{Рисунок 14 -- Результат в Swagger UI: Успешный POST /predict}
\end{figure}

\section{Организация хранения
результатов}\label{ux43eux440ux433ux430ux43dux438ux437ux430ux446ux438ux44f-ux445ux440ux430ux43dux435ux43dux438ux44f-ux440ux435ux437ux443ux43bux44cux442ux430ux442ux43eux432}

Создадим временное хранилище в оперативной памяти в файле storage.py
(Рисунок 15). Класс PredictionStorage использует словарь \_predictions,
где ключом является request\_id, а значением -- объект
PredictionResponse.

Метод save() сохраняет результат прогнозирования, get() находит
результат по идентификатору, а list() возвращает список всех сохранённых
прогнозов. Объект storage создаётся при запуске приложения и
используется другими компонентами сервиса.

Хранилище временное: после перезапуска приложения все сохранённые данные
теряются.

\begin{figure}
\centering
\includegraphics[width=5.92289in,height=3.85221in]{media/image14.png}
\caption{Рисунок 15 -- Реализация хранилища}
\end{figure}

\section{Реализация endpoint
/predictions/\{request\_id\}}\label{ux440ux435ux430ux43bux438ux437ux430ux446ux438ux44f-endpoint-predictionsrequest_id}

Уникальный request\_id необходим для того, чтобы каждый результат
прогнозирования можно было однозначно идентифицировать и получить с
помощью endpoint GET /predictions/\{request\_id\}. Реализуем данный
эндпоинт (Рисунок 16). И посмотрим результат в Swagger (Рисунок 17).

\begin{figure}
\centering
\includegraphics[width=6.73079in,height=1.80938in]{media/image15.png}
\caption{Рисунок 16 -- Реализация endpoint /predictions/\{request\_id\}}
\end{figure}

\begin{figure}
\centering
\includegraphics[width=6.36329in,height=4.50407in]{media/image16.png}
\caption{Рисунок 17 -- Результат в Swagger}
\end{figure}

Обратим внимание на ошибку 422. При передаче некорректных данных FastAPI
должен вернуть ошибку валидации (Рисунок 18).

\begin{figure}
\centering
\includegraphics[width=6.52419in,height=3.90806in]{media/image17.png}
\caption{Рисунок 18 -- Пример передачи некорректных данных: ошибка 422}
\end{figure}

Почему такой запрос не должен доходить до ML-модели?

Pydantic выполняет структурную валидацию до вызова функции predict:
проверяет типы, обязательные поля и ограничения Field. Например,
area\_ha должен быть больше нуля. При нарушении FastAPI возвращает 422,
поэтому инференс не запускается для некорректных входных данных.

Также попробуем ввести несуществующий request\_id (Рисунок 19).

\begin{figure}
\centering
\includegraphics[width=5.99333in,height=4.12364in]{media/image18.png}
\caption{Рисунок 19 -- Пример передачи несуществующих данных: ошибка
404}
\end{figure}

\section{Реализация endpoint
/predictions}\label{ux440ux435ux430ux43bux438ux437ux430ux446ux438ux44f-endpoint-predictions}

Теперь реализуем получение списка прогнозов (Рисунок 19). Endpoint
поддерживает два query-параметра: limit и risk\_level.

\begin{figure}
\centering
\includegraphics[width=6.61151in,height=2.26419in]{media/image19.png}
\caption{Рисунок 20 -- Реализация endpoint /predictions}
\end{figure}

Параметр limit задаёт максимальное количество возвращаемых прогнозов.
Для него устанавливаются ограничения от 1 до 100. Это позволяет избежать
некорректных запросов и получения слишком большого количества данных.

Параметр risk\_level используется для фильтрации прогнозов по уровню
риска: low, medium или high. Если параметр не указан, возвращаются все
сохранённые прогнозы.

Таким образом, можно выполнить запрос:

\begin{itemize}
\item
  /predictions?limit=2: получить максимум 2 прогноза (Рисунок 20);
\item
  /predictions?risk\_level=high: получить только прогнозы с высоким
  риском (Рисунок 21);
\item
  /predictions?limit=2\&risk\_level=high: получить до 2 прогнозов с
  высоким риском (Рисунок 22).
\end{itemize}

Получение данных выполняется через метод list() временного хранилища,
после чего при необходимости применяется фильтрация по уровню риска и
ограничение количества результатов.

\begin{figure}
\centering
\includegraphics[width=6.23084in,height=2.99215in]{media/image20.png}
\caption{Рисунок 21 -- Результат для /predictions?limit=2}
\end{figure}

\begin{figure}
\centering
\includegraphics[width=5.98327in,height=2.86077in]{media/image21.png}
\caption{Рисунок 22 -- Результат для /predictions?risk\_level=high}
\end{figure}

\begin{figure}
\centering
\includegraphics[width=6.06127in,height=2.94014in]{media/image22.png}
\caption{Рисунок 23 -- результат для
/predictions?limit=2\&risk\_level=high}
\end{figure}

При некорректном значении limit, например limit=1000, FastAPI возвращает
ошибку 422 Unprocessable Entity благодаря валидации query-параметра
(Рисунок 23).

\begin{figure}
\centering
\includegraphics[width=6.88233in,height=2.80639in]{media/image23.png}
\caption{Рисунок 24 -- Некорректный limit}
\end{figure}

\section{\texorpdfstring{Проверка бизнес ограничений
}{Проверка бизнес ограничений }}\label{ux43fux440ux43eux432ux435ux440ux43aux430-ux431ux438ux437ux43dux435ux441-ux43eux433ux440ux430ux43dux438ux447ux435ux43dux438ux439}

Pydantic проверяет формат данных, но некоторые ошибки невозможно
определить только по типу значения, иногда необходима дополнительная
проверка некоторых ограничений. Поэтому создадим множество допустимых
регионов (Рисунок 24). Также добавим валидацию в /predict (Рисунок 25).

\begin{figure}
\centering
\includegraphics[width=6.17406in,height=1.81127in]{media/image24.png}
\caption{Рисунок 25 -- Множество допустимых регионов}
\end{figure}

\begin{figure}
\centering
\includegraphics[width=6.69831in,height=1.31721in]{media/image25.png}
\caption{Рисунок 26 -- Валидация региона в /predict}
\end{figure}

Посмотрим на результат в Swagger при запросе с неизвестным регионом
(Рисунок 26).

\begin{figure}
\centering
\includegraphics[width=6.85873in,height=3.35272in]{media/image26.png}
\caption{Рисунок 27 -- Результат с неизвестным регионом: ошибка 400}
\end{figure}

\section{OpenAPI как
контракт}\label{openapi-ux43aux430ux43a-ux43aux43eux43dux442ux440ux430ux43aux442}

Откроем Swagger UI (Рисунок 27), в документации автоматически
отображаются:

\begin{itemize}
\item
  FarmRequest
\item
  PredictionResponse
\item
  типы данных
\item
  обязательные поля
\item
  ограничения
\item
  HTTP endpoints
\item
  HTTP methods
\end{itemize}

\begin{figure}
\centering
\includegraphics[width=6.12838in,height=4.92999in]{media/image27.png}
\caption{Рисунок 28 -- Скриншот Swagger UI}
\end{figure}

\section{Тестирование
API}\label{ux442ux435ux441ux442ux438ux440ux43eux432ux430ux43dux438ux435-api}

Последовательно выполним запросы и заполним таблицу тест-кейсов (Таблица
№2).

Таблица №2 -- Таблица тест-кейсов

\begin{longtable}[]{@{}
  >{\centering\arraybackslash}p{(\linewidth - 8\tabcolsep) * \real{0.0563}}
  >{\raggedright\arraybackslash}p{(\linewidth - 8\tabcolsep) * \real{0.2817}}
  >{\centering\arraybackslash}p{(\linewidth - 8\tabcolsep) * \real{0.1409}}
  >{\centering\arraybackslash}p{(\linewidth - 8\tabcolsep) * \real{0.1549}}
  >{\raggedright\arraybackslash}p{(\linewidth - 8\tabcolsep) * \real{0.3662}}@{}}
\toprule\noalign{}
\begin{minipage}[b]{\linewidth}\centering
\begin{quote}
\textbf{№}
\end{quote}
\end{minipage} & \begin{minipage}[b]{\linewidth}\centering
\begin{quote}
\textbf{Запрос}
\end{quote}
\end{minipage} & \begin{minipage}[b]{\linewidth}\centering
\begin{quote}
\textbf{Ожидаемый HTTP-код}
\end{quote}
\end{minipage} & \begin{minipage}[b]{\linewidth}\centering
\textbf{Полученный код}
\end{minipage} & \begin{minipage}[b]{\linewidth}\centering
\begin{quote}
\textbf{Результат}
\end{quote}
\end{minipage} \\
\midrule\noalign{}
\endhead
\bottomrule\noalign{}
\endlastfoot
\begin{minipage}[t]{\linewidth}\centering
\begin{quote}
1
\end{quote}
\end{minipage} & \begin{minipage}[t]{\linewidth}\raggedright
\begin{quote}
GET /health
\end{quote}
\end{minipage} & \begin{minipage}[t]{\linewidth}\centering
\begin{quote}
200
\end{quote}
\end{minipage} & \begin{minipage}[t]{\linewidth}\centering
\begin{quote}
200
\end{quote}
\end{minipage} & \begin{minipage}[t]{\linewidth}\raggedright
\begin{quote}
\{

"status": "ok"

\}
\end{quote}
\end{minipage} \\
\begin{minipage}[t]{\linewidth}\centering
\begin{quote}
2
\end{quote}
\end{minipage} & \begin{minipage}[t]{\linewidth}\raggedright
\begin{quote}
GET /model-info
\end{quote}
\end{minipage} & \begin{minipage}[t]{\linewidth}\centering
\begin{quote}
200
\end{quote}
\end{minipage} & \begin{minipage}[t]{\linewidth}\centering
\begin{quote}
200
\end{quote}
\end{minipage} & \begin{minipage}[t]{\linewidth}\raggedright
\begin{quote}
\{

"model\_name": "agro-risk-model",

"model\_version": "1.0",

"model\_type": "risk-scoring",

"status": "ready"

\}
\end{quote}
\end{minipage} \\
\begin{minipage}[t]{\linewidth}\centering
\begin{quote}
3
\end{quote}
\end{minipage} & \begin{minipage}[t]{\linewidth}\raggedright
\begin{quote}
корректный POST /predict
\end{quote}
\end{minipage} & \begin{minipage}[t]{\linewidth}\centering
\begin{quote}
200/201
\end{quote}
\end{minipage} & \begin{minipage}[t]{\linewidth}\centering
\begin{quote}
200/201
\end{quote}
\end{minipage} & \begin{minipage}[t]{\linewidth}\raggedright
\begin{quote}
\{

"request\_id": "ec8b3bcc-bccc-4794-865d-51511389cf78",

"farm\_id": "FARM-001",

"risk\_score": 0.9,

"risk\_level": "high",

"recommendation": "Высокий риск. Требуется ручное рассмотрение",

"model\_version": "1.0"

\}
\end{quote}
\end{minipage} \\
\begin{minipage}[t]{\linewidth}\centering
\begin{quote}
4
\end{quote}
\end{minipage} & \begin{minipage}[t]{\linewidth}\raggedright
\begin{quote}
area\_ha = -100
\end{quote}
\end{minipage} & \begin{minipage}[t]{\linewidth}\centering
\begin{quote}
422
\end{quote}
\end{minipage} & \begin{minipage}[t]{\linewidth}\centering
\begin{quote}
422
\end{quote}
\end{minipage} & \begin{minipage}[t]{\linewidth}\raggedright
\begin{quote}
\{

"detail": {[}

\{

"type": "greater\_than",

"loc": {[}

"body",

"area\_ha"

{]},

"msg": "Input should be greater than 0",

"input": -100,

"ctx": \{

"gt": 0

\}

\}

{]}

\}
\end{quote}
\end{minipage} \\
\begin{minipage}[t]{\linewidth}\centering
\begin{quote}
5
\end{quote}
\end{minipage} & \begin{minipage}[t]{\linewidth}\raggedright
\begin{quote}
неизвестный регион
\end{quote}
\end{minipage} & \begin{minipage}[t]{\linewidth}\centering
\begin{quote}
400
\end{quote}
\end{minipage} & \begin{minipage}[t]{\linewidth}\centering
\begin{quote}
400
\end{quote}
\end{minipage} & \begin{minipage}[t]{\linewidth}\raggedright
\begin{quote}
\{

"detail": "Unknown region: Moscow. Allowed regions:
{[}\textquotesingle Krasnodar\textquotesingle,
\textquotesingle Rostov\textquotesingle,
\textquotesingle Stavropol\textquotesingle{]}"

\}
\end{quote}
\end{minipage} \\
\begin{minipage}[t]{\linewidth}\centering
\begin{quote}
6
\end{quote}
\end{minipage} & \begin{minipage}[t]{\linewidth}\raggedright
\begin{quote}
существующий request\_id
\end{quote}
\end{minipage} & \begin{minipage}[t]{\linewidth}\centering
\begin{quote}
200
\end{quote}
\end{minipage} & \begin{minipage}[t]{\linewidth}\centering
\begin{quote}
200
\end{quote}
\end{minipage} & \begin{minipage}[t]{\linewidth}\raggedright
\begin{quote}
\{

"request\_id": "ec8b3bcc-bccc-4794-865d-51511389cf78",

"farm\_id": "FARM-001",

"risk\_score": 0.9,

"risk\_level": "high",

"recommendation": "Высокий риск. Требуется ручное рассмотрение",

"model\_version": "1.0"

\}
\end{quote}
\end{minipage} \\
\end{longtable}

Продолжение Таблицы №2

\begin{longtable}[]{@{}
  >{\centering\arraybackslash}p{(\linewidth - 8\tabcolsep) * \real{0.0563}}
  >{\raggedright\arraybackslash}p{(\linewidth - 8\tabcolsep) * \real{0.2817}}
  >{\centering\arraybackslash}p{(\linewidth - 8\tabcolsep) * \real{0.1409}}
  >{\centering\arraybackslash}p{(\linewidth - 8\tabcolsep) * \real{0.1549}}
  >{\raggedright\arraybackslash}p{(\linewidth - 8\tabcolsep) * \real{0.3662}}@{}}
\toprule\noalign{}
\begin{minipage}[b]{\linewidth}\centering
\begin{quote}
\textbf{№}
\end{quote}
\end{minipage} & \begin{minipage}[b]{\linewidth}\raggedright
\begin{quote}
\textbf{Запрос}
\end{quote}
\end{minipage} & \begin{minipage}[b]{\linewidth}\centering
\begin{quote}
\textbf{Ожидаемый HTTP-код}
\end{quote}
\end{minipage} & \begin{minipage}[b]{\linewidth}\centering
\begin{quote}
\textbf{Полученный код}
\end{quote}
\end{minipage} & \begin{minipage}[b]{\linewidth}\raggedright
\begin{quote}
\textbf{Результат}
\end{quote}
\end{minipage} \\
\midrule\noalign{}
\endhead
\bottomrule\noalign{}
\endlastfoot
\begin{minipage}[t]{\linewidth}\centering
\begin{quote}
7
\end{quote}
\end{minipage} & \begin{minipage}[t]{\linewidth}\raggedright
\begin{quote}
неизвестный request\_id
\end{quote}
\end{minipage} & \begin{minipage}[t]{\linewidth}\centering
\begin{quote}
404
\end{quote}
\end{minipage} & \begin{minipage}[t]{\linewidth}\centering
\begin{quote}
404
\end{quote}
\end{minipage} & \begin{minipage}[t]{\linewidth}\raggedright
\begin{quote}
\{

"detail": "Prediction not found"

\}
\end{quote}
\end{minipage} \\
\begin{minipage}[t]{\linewidth}\centering
\begin{quote}
8
\end{quote}
\end{minipage} & \begin{minipage}[t]{\linewidth}\raggedright
\begin{quote}
limit = 2
\end{quote}
\end{minipage} & \begin{minipage}[t]{\linewidth}\centering
\begin{quote}
200
\end{quote}
\end{minipage} & \begin{minipage}[t]{\linewidth}\centering
\begin{quote}
200
\end{quote}
\end{minipage} & \begin{minipage}[t]{\linewidth}\raggedright
\begin{quote}
{[}

\{

"request\_id": "80f7160b-9ea5-41f8-8759-51e04143f608",

"farm\_id": "FARM-001",

"risk\_score": 0.9,

"risk\_level": "high",

"recommendation": "Высокий риск. Требуется ручное рассмотрение",

"model\_version": "1.0"

\},

\{

"request\_id": "11a3831e-5a02-45e3-b0e8-79c69d9d0aaa",

"farm\_id": "FARM-001",

"risk\_score": 0.7,

"risk\_level": "high",

"recommendation": "Высокий риск. Требуется ручное рассмотрение",

"model\_version": "1.0"

\}

{]}
\end{quote}
\end{minipage} \\
\begin{minipage}[t]{\linewidth}\centering
\begin{quote}
9
\end{quote}
\end{minipage} & \begin{minipage}[t]{\linewidth}\raggedright
\begin{quote}
risk\_level = high
\end{quote}
\end{minipage} & \begin{minipage}[t]{\linewidth}\centering
\begin{quote}
200
\end{quote}
\end{minipage} & \begin{minipage}[t]{\linewidth}\centering
\begin{quote}
200
\end{quote}
\end{minipage} & \begin{minipage}[t]{\linewidth}\raggedright
\begin{quote}
{[}

\{

"request\_id": "80f7160b-9ea5-41f8-8759-51e04143f608",

"farm\_id": "FARM-001",

"risk\_score": 0.9,

"risk\_level": "high",

"recommendation": "Высокий риск. Требуется ручное рассмотрение",

"model\_version": "1.0"

\},

\{

"request\_id": "11a3831e-5a02-45e3-b0e8-79c69d9d0aaa",

"farm\_id": "FARM-001",

"risk\_score": 0.7,

"risk\_level": "high",

"recommendation": "Высокий риск. Требуется ручное рассмотрение",

"model\_version": "1.0"

\}

{]}
\end{quote}
\end{minipage} \\
\begin{minipage}[t]{\linewidth}\centering
\begin{quote}
10
\end{quote}
\end{minipage} & \begin{minipage}[t]{\linewidth}\raggedright
\begin{quote}
limit = -5
\end{quote}
\end{minipage} & \begin{minipage}[t]{\linewidth}\centering
\begin{quote}
422
\end{quote}
\end{minipage} & \begin{minipage}[t]{\linewidth}\centering
\begin{quote}
422
\end{quote}
\end{minipage} & \begin{minipage}[t]{\linewidth}\raggedright
Please correct the following validation errors and try again.

\begin{itemize}
\item
  For \textquotesingle limit\textquotesingle: Value must be greater than
  or equal to 1.
\end{itemize}
\end{minipage} \\
\end{longtable}

\section{Анализ API как
контракта}\label{ux430ux43dux430ux43bux438ux437-api-ux43aux430ux43a-ux43aux43eux43dux442ux440ux430ux43aux442ux430}

Для каждого endpoint заполним таблицу (Таблица №5).

Таблица №5 -- Анализ API-контракта

\begin{longtable}[]{@{}
  >{\raggedright\arraybackslash}p{(\linewidth - 10\tabcolsep) * \real{0.1673}}
  >{\centering\arraybackslash}p{(\linewidth - 10\tabcolsep) * \real{0.0935}}
  >{\raggedright\arraybackslash}p{(\linewidth - 10\tabcolsep) * \real{0.1633}}
  >{\raggedright\arraybackslash}p{(\linewidth - 10\tabcolsep) * \real{0.2791}}
  >{\centering\arraybackslash}p{(\linewidth - 10\tabcolsep) * \real{0.0911}}
  >{\raggedright\arraybackslash}p{(\linewidth - 10\tabcolsep) * \real{0.2057}}@{}}
\toprule\noalign{}
\begin{minipage}[b]{\linewidth}\centering
\textbf{Endpoint}
\end{minipage} & \begin{minipage}[b]{\linewidth}\centering
\textbf{Method}
\end{minipage} & \begin{minipage}[b]{\linewidth}\centering
\textbf{Input}
\end{minipage} & \begin{minipage}[b]{\linewidth}\centering
\textbf{Output}
\end{minipage} & \begin{minipage}[b]{\linewidth}\centering
\textbf{Success}
\end{minipage} & \begin{minipage}[b]{\linewidth}\centering
\textbf{Ошибки}
\end{minipage} \\
\midrule\noalign{}
\endhead
\bottomrule\noalign{}
\endlastfoot
/health & GET & --- & HealthResponse -- поле status со значением "ok" &
200 & 500 при необработанном внутреннем сбое \\
/model-info & GET & --- & ModelInfoResponse -- model\_name,
model\_version, model\_type, status & 200 & 500 при необработанном
внутреннем сбое \\
/predict & POST & FarmRequest & PredictionResponse -- request\_id,
farm\_id, risk\_score, risk\_level, recommendation, model\_version & 200
& 400, 422, 503 \\
/predictions/\{id\} & GET & request\_id & PredictionResponse & 200 &
404 \\
/predictions & GET & query params: limit, risk\_level &
list{[}PredictionResponse{]} & 200 & 400, 422 \\
\end{longtable}

Статус POST /predict в данном репозитории -- именно 200 OK, поскольку
так задан декоратор маршрута. Код 500 обозначает потенциальный
необработанный сбой, а не отдельную явно реализованную ветку. Неверный
risk\_level даёт 400, так как проверяется вручную в обработчике.

Проведённый анализ API-контракта показывает, что сервис использует
REST-подход и HTTP-методы GET и POST. Для передачи данных
прогнозирования используется JSON, а структура входных и выходных данных
определяется Pydantic-моделями. Для получения конкретного прогноза
используется path-параметр request\_id, а для фильтрации списка
прогнозов -- query-параметры limit и risk\_level.

Корректность входных данных проверяется до выполнения бизнес-логики. При
нарушении структурных ограничений возвращается HTTP 422, при неизвестном
регионе -- HTTP 400, а при отсутствии запрошенного прогноза -- HTTP 404.
Успешное выполнение запросов возвращает HTTP 200.

\section{Разделение приложения на модули: архитектура
приложения}\label{ux440ux430ux437ux434ux435ux43bux435ux43dux438ux435-ux43fux440ux438ux43bux43eux436ux435ux43dux438ux44f-ux43dux430-ux43cux43eux434ux443ux43bux438-ux430ux440ux445ux438ux442ux435ux43aux442ux443ux440ux430-ux43fux440ux438ux43bux43eux436ux435ux43dux438ux44f}

Разработан учебный REST API агроскоринга на FastAPI. Его потребители --
CRM банка, веб-портал, мобильное приложение, внутренние банковские
сервисы и BI-система. Клиенты отправляют сведения о хозяйстве в JSON по
HTTP и получают числовую оценку риска, категорию риска и рекомендацию.
FastAPI автоматически публикует контракт OpenAPI по адресу /openapi.json
и интерфейс Swagger UI по адресу /docs.

В процесесе выполнения практической работы код был приведён к модульной
структуре:

\begin{itemize}
\item
  agro\_api/main.py: FastAPI-приложение и HTTP endpoints;
\item
  agro\_api/schemas.py: Pydantic-модели;
\item
  agro\_api/model.py: детерминированная функция расчёта risk\_score;
\item
  agro\_api/services.py: бизнес-логика и postprocessing;
\item
  agro\_api/storage.py: временное in-memory хранилище;
\item
  tests/test\_api.py: автоматические тесты API.
\end{itemize}

Прогноз выполняется синхронно в рамках POST /predict. Сначала данные
проходят валидацию Pydantic, затем проверяются регион и готовность
модели. После расчёта и постобработки результат получает request\_id,
сохраняется и возвращается клиенту в формате PredictionResponse. Такая
декомпозиция уменьшает связанность и упрощает дальнейшую замену правил
расчёта на реальную ML-модель и временного хранилища на БД.

\section{ВЫВОД}\label{ux432ux44bux432ux43eux434}

Компонентная схема архитектуры (Рисунок 29) показывает клиентов,
HTTP-слой FastAPI, Pydantic-схемы, сервисную логику, функцию расчёта
риска и хранилище. REST API служит стабильной границей между внешними
банковскими системами и расчётом риска: принимает и проверяет данные,
запускает инференс, документирует формат результата и сообщает об
ошибках через HTTP-коды.

\begin{figure}
\centering
\includegraphics[width=6.48854in,height=2.89245in]{media/image28.png}
\caption{Рисунок 29 --Компонентная схема архитектуры}
\end{figure}

Текущая модель является учебным правилом, а хранилище работает только в
памяти процесса: результаты исчезают после перезапуска. В реализации
storage.list() записи идут в порядке добавления, поэтому limit выбирает
первые сохранённые прогнозы, а не самые новые. Для промышленной системы
понадобятся постоянное хранилище, защита доступа, аудит и обученная
модель с контролем качества.

Таким образом, в данной практической работе реализован и проверен REST
API для AI-сервиса оценки риска. FastAPI предоставляет HTTP-интерфейс и
OpenAPI-контракт, Pydantic обеспечивает структурную валидацию,
бизнес-валидация защищает доменные ограничения, а request\_id позволяет
получать сохранённый результат. Рефакторинг разделяет HTTP-слой, схемы,
модель, сервисную логику и хранилище. REST API в AI-системе выступает
контрактом между клиентом и инференсным сервисом.

\section{Приложение А: Полный код
приложения}\label{ux43fux440ux438ux43bux43eux436ux435ux43dux438ux435-ux430-ux43fux43eux43bux43dux44bux439-ux43aux43eux434-ux43fux440ux438ux43bux43eux436ux435ux43dux438ux44f}

Полный код приложения находится в репозитории:

\url{https://github.com/Shynuaa/MINS_labs}

\end{document}
